% Options for packages loaded elsewhere
\PassOptionsToPackage{unicode}{hyperref}
\PassOptionsToPackage{hyphens}{url}
\PassOptionsToPackage{dvipsnames,svgnames,x11names}{xcolor}
%
\documentclass[
  11pt,
  a4paper,
]{report}
\usepackage{amsmath,amssymb}
\usepackage{iftex}
\ifPDFTeX
  \usepackage[T1]{fontenc}
  \usepackage[utf8]{inputenc}
  \usepackage{textcomp} % provide euro and other symbols
\else % if luatex or xetex
  \usepackage{unicode-math} % this also loads fontspec
  \defaultfontfeatures{Scale=MatchLowercase}
  \defaultfontfeatures[\rmfamily]{Ligatures=TeX,Scale=1}
\fi
\usepackage{lmodern}
\ifPDFTeX\else
  % xetex/luatex font selection
\fi
% Use upquote if available, for straight quotes in verbatim environments
\IfFileExists{upquote.sty}{\usepackage{upquote}}{}
\IfFileExists{microtype.sty}{% use microtype if available
  \usepackage[]{microtype}
  \UseMicrotypeSet[protrusion]{basicmath} % disable protrusion for tt fonts
}{}
\makeatletter
\@ifundefined{KOMAClassName}{% if non-KOMA class
  \IfFileExists{parskip.sty}{%
    \usepackage{parskip}
  }{% else
    \setlength{\parindent}{0pt}
    \setlength{\parskip}{6pt plus 2pt minus 1pt}}
}{% if KOMA class
  \KOMAoptions{parskip=half}}
\makeatother
\usepackage{xcolor}
\usepackage[margin=25mm]{geometry}
\usepackage{longtable,booktabs,array}
\usepackage{calc} % for calculating minipage widths
% Correct order of tables after \paragraph or \subparagraph
\usepackage{etoolbox}
\makeatletter
\patchcmd\longtable{\par}{\if@noskipsec\mbox{}\fi\par}{}{}
\makeatother
% Allow footnotes in longtable head/foot
\IfFileExists{footnotehyper.sty}{\usepackage{footnotehyper}}{\usepackage{footnote}}
\makesavenoteenv{longtable}
\usepackage{graphicx}
\makeatletter
\def\maxwidth{\ifdim\Gin@nat@width>\linewidth\linewidth\else\Gin@nat@width\fi}
\def\maxheight{\ifdim\Gin@nat@height>\textheight\textheight\else\Gin@nat@height\fi}
\makeatother
% Scale images if necessary, so that they will not overflow the page
% margins by default, and it is still possible to overwrite the defaults
% using explicit options in \includegraphics[width, height, ...]{}
\setkeys{Gin}{width=\maxwidth,height=\maxheight,keepaspectratio}
% Set default figure placement to htbp
\makeatletter
\def\fps@figure{htbp}
\makeatother
\setlength{\emergencystretch}{3em} % prevent overfull lines
\providecommand{\tightlist}{%
  \setlength{\itemsep}{0pt}\setlength{\parskip}{0pt}}
\setcounter{secnumdepth}{5}
\ifLuaTeX
  \usepackage{selnolig}  % disable illegal ligatures
\fi
\IfFileExists{bookmark.sty}{\usepackage{bookmark}}{\usepackage{hyperref}}
\IfFileExists{xurl.sty}{\usepackage{xurl}}{} % add URL line breaks if available
\urlstyle{same}
\hypersetup{
  pdftitle={VIGILANT: Change-Aware Software Traceability and Evidence-Grounded Investigation},
  colorlinks=true,
  linkcolor={Maroon},
  filecolor={Maroon},
  citecolor={Blue},
  urlcolor={Blue},
  pdfcreator={LaTeX via pandoc}}

\title{VIGILANT: Change-Aware Software Traceability and
Evidence-Grounded Investigation}
\usepackage{etoolbox}
\makeatletter
\providecommand{\subtitle}[1]{% add subtitle to \maketitle
  \apptocmd{\@title}{\par {\large #1 \par}}{}{}
}
\makeatother
\subtitle{M.Tech technical thesis draft for review}
\author{}
\date{8 October 2026}

\begin{document}
\maketitle

{
\hypersetup{linkcolor=}
\setcounter{tocdepth}{0}
\tableofcontents
}
\hypertarget{introduction}{%
\chapter{Introduction}\label{introduction}}

Software maintenance involves relationships that cross artifact
boundaries. A requirement can describe behavior implemented in a backend
handler, invoked by a generated frontend client and exercised by tests.
A change to one artifact may leave the others unchanged. A useful
investigation must identify which versions differ, locate plausible
dependencies and provide enough source evidence for a reviewer to assess
whether the change matters.

VIGILANT implements this investigation as a reproducible pipeline over
Git repositories. It combines canonical artifacts, lexical and
pretrained embedding retrieval, a structural knowledge graph, bounded
impact analysis, retrieval-supported model calls and conservative
consistency rules. The intended output is an inspectable report
containing changed artifacts, affected candidates, trace candidates,
evidence, tool outcomes and a separate model decision. Evidence retains
the repository, commit, path, line range and content hash required to
reopen its origin.

This thesis distinguishes an engineering achievement from an
effectiveness claim. The implementation executes end to end, and its
completed studies contain actual rankings and model generations.
However, the available results do not demonstrate that the full
local-model system classifies consistency better than an
always-UNCERTAIN comparator. Retrieval effects also vary by dataset and
metric. Reporting these outcomes is necessary to make the contribution
assessable and reproducible.

The document is a technical draft. Author, supervisor, institution,
prescribed front matter and final university formatting require
confirmation. Independent semantic review, external plagiarism checks,
browser validation and Windows execution remain open. No publication
submission or venue-indexing claim accompanies this draft.

\hypertarget{background}{%
\chapter{Background}\label{background}}

Traceability link recovery ranks or classifies relationships between
source and target artifacts. Ranking is useful when a human examines a
short candidate list; binary classification answers whether a link
should be accepted. These are different decision problems. A high-ranked
item can still be incorrect, and a ranking score is not automatically a
calibrated probability of implementation.

Lexical retrieval uses token overlap and weighting. BM25 provides a
probabilistic retrieval framework {[}1{]}. Sentence representations
support cosine-based comparison of separately encoded texts {[}3{]}.
VIGILANT uses a pinned MiniLM encoder {[}10{]} and combines BM25 and
embedding ranks through reciprocal rank fusion {[}2{]}. The implemented
fusion uses a constant of 60 and two retrieval arms. This setting is a
configuration choice, not proof that fusion must improve every dataset.

A software knowledge graph records typed artifact relationships. In this
project, parsed calls, imports, tests, APIs and client relationships
supply structural context. A graph path is a reason to investigate an
artifact, not evidence that the artifact is defective. Likewise, a
recovered similarity candidate is not a proven IMPLEMENTS relation.

Retrieval-augmented generation supplies external context to a language
model {[}4{]}. VIGILANT uses this idea for candidate review and
consistency investigation. The original RAG formulation, the
traceability literature and the implemented pipeline have different
training and evaluation settings. The present experiments evaluate the
implemented method directly rather than claiming a reproduction of those
systems.

\hypertarget{literature-review}{%
\chapter{Literature Review}\label{literature-review}}

The selected references cover retrieval, software traceability,
graph-supported retrieval and tool-using agents. This is a targeted
review based on primary metadata, abstracts and model documentation
checked on 7 October 2026. It is not an exhaustive systematic review.
The separate survey draft should expand full-text screening before any
comprehensive literature claim.

LiSSA investigates RAG-based linking across requirements, code and
architecture artifacts {[}5{]}. TransArC uses an intermediate
architecture model to connect documentation and code {[}6{]}. Earlier
ARDoCo-related work uses trace links to identify documentation/model
inconsistencies {[}7{]}. These studies demonstrate that automated
linking and inconsistency analysis already have substantial prior work;
VIGILANT does not claim to be the first system to combine language
models with traceability.

GraphRAG addresses corpus-level summarization through an entity graph
and community summaries {[}8{]}. VIGILANT instead constructs a parsed
software graph and performs bounded local traversal. SWE-agent studies
software-engineering agents and their computer interfaces {[}9{]}.
VIGILANT's current investigator follows a scripted sequence of tools.
These relationships motivate design and evaluation questions, but their
reported effectiveness is not numerically comparable with the present
traceability and mutation experiments.

The experimental comparison in this thesis uses only VIGILANT's own
rankings and predictions under its recorded populations. External paper
scores are not copied into the results as if they were reproduced
baselines. The reference-verification record identifies the sources and
access limitations.

\hypertarget{research-gap}{%
\chapter{Research Gap}\label{research-gap}}

The project targets the practical integration of version-aware artifact
extraction, trace candidate recovery, structural context and inspectable
model-supported investigation. Static trace rankings alone do not
describe how a particular Git change propagates through a frontend,
backend and tests. Conversely, a model's fluent account of a change is
difficult to evaluate unless its evidence is tied to actual artifact
versions.

The proposed contribution is an integrated experimental artifact and an
explicit evaluation of its boundaries. It includes a controlled change
benchmark with reproducible commits, a frozen multi-dataset retrieval
study, complete model-call records and negative ablation results. This
gap statement is a project motivation, not a verified claim that no
prior system addresses software evolution or evidence provenance.

The most important unresolved research gap inside the current project is
independent validation. Authored mutations and source-faithful citations
are insufficient to establish reliable semantic judgments on naturally
occurring changes. That limitation constrains the conclusion even when
the software executes successfully.

\hypertarget{problem-statement}{%
\chapter{Problem Statement}\label{problem-statement}}

Given a repository, a baseline Git commit and a changed commit,
determine which canonical artifacts changed, which related artifacts
deserve investigation, and which traceability or consistency findings
can be supported by recorded evidence. The system must retain
uncertainty when the available context does not justify a conclusion.

There are three related outputs. First, retrieval methods produce
ordered trace candidates for a requirement. Second, graph traversal
produces impact candidates with paths explaining their selection. Third,
rules and model calls produce findings with source evidence. A candidate
list, a structural path and a consistency decision must remain
distinguishable in storage, the interface and the evaluation.

The system must not treat a missing test update as automatic proof of a
defect, a lexical match as confirmed implementation, or a loading error
as a generated model prediction. Reproducibility requires the recorded
result to identify its dataset, input, method, model, configuration,
output, label scope, metric and source record.

\hypertarget{objectives}{%
\chapter{Objectives}\label{objectives}}

The implementation objectives are to import Git snapshots, extract
supported artifact types, assign stable identities and content versions,
recover bounded trace candidates, persist structural graphs and execute
a complete evidence-backed change investigation. Runtime failures and
uncertain decisions must remain visible.

The evaluation objectives are to compare TF-IDF, BM25, MiniLM, hybrid
retrieval, graph reranking and a frozen RAG reranker; execute controlled
software-change scenarios; and compare five local-model investigation
variants. All comparisons use retained raw outputs and explicit
ground-truth scope.

The delivery objectives are to preserve one project root with code,
data, results, graphs, scripts and resume instructions, and to prepare
reviewable academic drafts and a demonstration. Final acceptance by the
university or a publication venue is outside the evidence available to
this project run.

\hypertarget{research-questions}{%
\chapter{Research Questions}\label{research-questions}}

\textbf{RQ1.} How do lexical, embedding, hybrid, graph and RAG methods
compare on the selected linked-query traceability populations?

\textbf{RQ2.} Does adding retrieval, graph context, bounded
investigation and temporal evidence improve local-model agreement with
the authored consistency labels?

\textbf{RQ3.} Can the implemented pipeline recover real Git changes,
investigate related artifacts and retain inspectable evidence and tool
outcomes?

\textbf{RQ4.} What aspects of grounding and reliability are established
by source/schema audits, and what still requires independent review?

RQ1 concerns recorded ranking quality. RQ2 concerns the local-model
pilot, not every possible model or prompt. RQ3 concerns executed
engineering behavior. RQ4 prevents structural validation from being
reported as semantic accuracy.

\hypertarget{vigilant-architecture}{%
\chapter{VIGILANT Architecture}\label{vigilant-architecture}}

The ingestion layer resolves commits and extracts canonical artifacts
from text snapshots. The persistence layer stores projects,
repositories, commits, artifact versions, trace links, graph nodes and
edges, evidence, findings, analysis runs and experiments. Retrieval
indexes target text; optional semantic retrieval uses the pinned encoder
and a content/model-aware cache.

The graph layer builds a directed multigraph at each commit. The
investigation composes base and head graphs to traverse around changed
artifacts. Evidence retrieval combines ranked text with graph and change
context. A verifier checks quotes, ranges and hashes before the provider
receives selected excerpts. Rules produce conservative findings
separately from the model's structured decision.

The local API exposes stored records and queues one bounded analysis
worker. The React frontend reads the API for dashboard, repository,
artifact, traceability, graph, change, impact, consistency, evidence,
agent, experiment and report views. The production build exists; browser
interaction validation remains blocked by the unavailable browser
runtime.

The investigation has named roles but a predetermined order. The saved
demonstration contains commit resolution, snapshot import, change
detection, two graph builds, impact analysis, trace recovery, evidence
retrieval, deterministic checks, quote verification, a model call and
report persistence. This is an executed tool workflow, not evidence of
autonomous plan selection.

\hypertarget{methodology}{%
\chapter{Methodology}\label{methodology}}

\hypertarget{canonical-identity-and-versions}{%
\section{Canonical identity and
versions}\label{canonical-identity-and-versions}}

An artifact ID hashes the project, repository, relative path, type and
symbol. Content has a separate hash and the artifact retains its commit
and line range. This separation supports comparison of logical artifacts
across snapshots while preserving the exact text used as evidence. Path
or symbol renames require explicit change matching; stable identity
alone does not solve every rename.

\hypertarget{retrieval-and-candidate-bounds}{%
\section{Retrieval and candidate
bounds}\label{retrieval-and-candidate-bounds}}

Token processing splits camel case, normalizes case and removes selected
requirement identifier patterns. TF-IDF uses sublinear term frequency.
BM25 and embedding arms rank the same target population. The encoder
limits text to up to eight 180-word chunks, averages normalized chunk
embeddings and normalizes the resulting vector. This can omit later code
in a long artifact.

Hybrid scores sum reciprocal ranks from BM25 and embeddings. The
implementation's automatic trace recovery returns at most three
candidates per requirement, with a 0.12 threshold in the lexical
fallback. These candidates retain RELATED\_TO and a proposed relation
with review-required status. No calibrated confidence is asserted.

\hypertarget{graph-context}{%
\section{Graph context}\label{graph-context}}

Structural relationships include CALLS, IMPORTS, TESTS, EXPOSES\_API,
CONSUMES\_API, DEPENDS\_ON and commit MODIFIES edges. Unique static
matches are preferred and ambiguous or dynamic relationships can be
unresolved. The change workflow uses bounded, typically two-hop
traversal and excludes broad reverse-containment propagation.

The retrieval graph experiment uses Java file relationships without gold
trace edges. Its exploratory score combines normalized hybrid relevance
with mean adjacent-file relevance using weights 0.85 and 0.15. This
differs from the operational impact graph and from Microsoft GraphRAG.
RTMX graph reranking was not run under this Java-file protocol.

\hypertarget{frozen-rag-reranking}{%
\section{Frozen RAG reranking}\label{frozen-rag-reranking}}

For each linked query, the model receives a requirement prefix of at
most 900 characters and five source/test prefixes of at most 1100
characters. Qwen assigns SUPPORTED, UNSUPPORTED or UNCERTAIN to each
candidate. SUPPORTED items move to the front while all other items
retain their relative order. The rest of the full hybrid ranking remains
available for scoring.

The model cannot introduce candidates outside the first five. Recall@5
and Recall@10 therefore cannot improve under this experiment. MAP,
reciprocal rank and nDCG can change. Prompt, evidence, schema, provider
source, model/decode identity and prepared-data fingerprints must match
for a cached generation to be reused.

\hypertarget{evidence-and-conservative-decisions}{%
\section{Evidence and conservative
decisions}\label{evidence-and-conservative-decisions}}

Evidence records contain exact text, path, line range, commit, artifact
identity and hash. The verifier compares each quote with the specified
artifact version. Model citations must refer to supplied IDs. Invalid
generated responses are retained and handled conservatively; loading
failures without output are excluded from effectiveness scoring.

The consistency decision schema requires coherent labels, reason,
uncertainty, affected IDs, citation IDs and confidence within {[}0,1{]}.
Non-UNCERTAIN conclusions require evidence. These checks enforce format
and provenance. They do not independently decide whether the excerpt
entails the conclusion.

\hypertarget{dataset-and-benchmark}{%
\chapter{Dataset and Benchmark}\label{dataset-and-benchmark}}

\begin{longtable}[]{@{}llllll@{}}
\toprule\noalign{}
Dataset & Split & Requirements & Linked & Targets & Links \\
\midrule\noalign{}
\endhead
\bottomrule\noalign{}
\endlastfoot
SMOS & train & 67 & 67 & 100 & 1044 \\
iTrust & validation & 131 & 105 & 226 & 286 \\
eTour\_en & test & 58 & 57 & 116 & 308 \\
RTMX-module-valid & external\_test & 22 & 22 & 7 & 22 \\
\end{longtable}

Only requirements with at least one resolved gold link enter the ranking
evaluation. This excludes 26 iTrust requirements and one eTour
requirement. The policy is identical across reported retrieval methods
but does not measure rejection of unlinked requirements.

SMOS, iTrust and eTour data were acquired from the pinned LiSSA
replication package. Their train, validation and test roles apply to the
learned-ranker study. RTMX is an external module-level subset with 22
resolved requirements and seven existing target modules. Missing modules
and unresolved function labels prevent interpreting it as a comparable
24-function evaluation. Returning all seven candidates makes Recall@10
equal one regardless of ranking quality.

VIGILANT-ChangeBench stores repository, baseline commit, changed commit,
change type, changed artifacts, expected affected artifacts and expected
consistency status. Its twelve scenarios include requirement
change/addition, code change/deletion, test change/deletion, frontend
change, API change, backend change, consumer change and
implementation/requirement test drift. The commits are real and
reproducible; the labels are authored and partial.

\begin{longtable}[]{@{}llll@{}}
\toprule\noalign{}
Case & Expected status & Known hits & Known misses \\
\midrule\noalign{}
\endhead
\bottomrule\noalign{}
\endlastfoot
01-requirement-changed & POTENTIAL INCONSISTENCY & 1 & 0 \\
02-requirement-added & UNCERTAIN & 0 & 0 \\
03-code-changed & POTENTIAL INCONSISTENCY & 1 & 0 \\
04-code-deleted & BROKEN TRACE & 1 & 0 \\
05-test-changed & UNCERTAIN & 1 & 0 \\
06-test-deleted & MISSING TEST & 1 & 0 \\
07-frontend-changed & UNCERTAIN & 0 & 0 \\
08-api-changed & POTENTIAL INCONSISTENCY & 5 & 0 \\
09-backend-changed & POTENTIAL INCONSISTENCY & 1 & 0 \\
10-consumer-changed & BROKEN TRACE & 1 & 4 \\
11-implementation-test-drift & POTENTIAL INCONSISTENCY & 1 & 0 \\
12-requirement-test-drift & POTENTIAL INCONSISTENCY & 1 & 0 \\
\end{longtable}

The benchmark seeds declared baseline traces in relevant scenarios. It
therefore checks maintenance of supplied relationships as well as
structural changes, rather than independently evaluating initial trace
discovery. Predictions outside a partial expected set remain unjudged.
SAFA is reported separately as a versioned artifact-tree comparison.

\hypertarget{implementation}{%
\chapter{Implementation}\label{implementation}}

Python modules separate ingestion, retrieval, traceability, graph
construction, change detection, evidence, model providers, consistency
and workflow orchestration. FastAPI exposes the local application and
SQLAlchemy stores records. The frontend uses React and Vite. Python AST
and language-specific parsing support the stated source languages, with
failures retained instead of silently disappearing from reports.

The local model provider uses Transformers on CPU with four inference
threads, greedy decoding, bfloat16 weights, an input limit of 3072
tokens and at most 220 generated tokens. Schema-constrained generation
restricts output structure. The pinned Qwen model revision is
\texttt{989aa7980e4cf806f80c7fef2b1adb7bc71aa306}. The MiniLM revision
is \texttt{1110a243fdf4706b3f48f1d95db1a4f5529b4d41}. Cloud-compatible,
Gemini, Anthropic and Ollama adapters exist but are not evaluated here.

Efficiency controls include embedding caches, bounded chunks, limited
candidate counts, graph-hop limits and exact reuse of generated
observations. The continuation ledger skips completed stages. These
controls reduce repeated work; no controlled universal speedup result is
claimed. Raw calls and rankings remain immutable while derived reports
can be regenerated.

The project also includes installation, restoration, checkpoint and
packaging scripts. SQLite backups preserve consistent snapshots. Archive
creation stages and validates ZIPs before installing them. This was
necessary because temporary runtimes repeatedly reset during evaluation.
Recovered calls retain their original timings and provenance rather than
being regenerated to obtain more favorable outcomes.

\hypertarget{experimental-setup}{%
\chapter{Experimental Setup}\label{experimental-setup}}

The frozen linked-query study includes 251 actual local generations.
Four completed dataset pointers identify immutable runs. Independent
audits compare prompt/evidence fingerprints, raw and parsed output,
schema constraints, full rankings and recomputed metrics. They verified
all observations and 23 method rows containing 1484 per-query metric
rows. A continuation check with zero new calls skipped every completed
dataset.

Precision@10, recall@10 and F1@10 are micro-averaged over predicted
requirement-target pairs. Macro Recall@k averages each query's fraction
of recovered gold links. These quantities can differ. MRR averages the
reciprocal rank of the first relevant target. MAP averages per-query
average precision across the full ranking. nDCG uses a cutoff of ten
with binary relevance. The complete raw and unrounded values remain
available.

For query q with gold set G and a full ranking, average precision sums
precision at every relevant rank and divides by the size of G. This
differs from measuring only the first five model-reviewed candidates. A
reranking operation can improve one order-sensitive metric and reduce
another without changing cutoff recall.

The learned ranker uses TF-IDF, normalized BM25, embedding and
token-Jaccard features. It is trained on SMOS, checked on iTrust and
evaluated on eTour. Its retention rule requires validation MAP to exceed
the strongest retrieval baseline. No test-based retention decision is
used. The fixed graph weight is exploratory rather than preregistered.

The local-model ablation uses the same twelve controlled cases across
five configurations, producing 60 observations. E is LLM-only; F adds
hybrid RAG; G adds graph context; H adds bounded investigation; I adds
temporal context and rules. H is scripted, and I changes multiple
context elements together, limiting causal attribution.

\hypertarget{results}{%
\chapter{Results}\label{results}}

\hypertarget{traceability}{%
\section{Traceability}\label{traceability}}

\begin{longtable}[]{@{}llllll@{}}
\toprule\noalign{}
Dataset & Method & MAP & MRR & nDCG@10 & Micro F1@10 \\
\midrule\noalign{}
\endhead
\bottomrule\noalign{}
\endlastfoot
SMOS & TF-IDF & 0.2924 & 0.6999 & 0.3530 & 0.2275 \\
SMOS & BM25 & 0.2979 & 0.6535 & 0.3334 & 0.2275 \\
iTrust & TF-IDF & 0.3160 & 0.4439 & 0.3860 & 0.2081 \\
iTrust & BM25 & 0.3787 & 0.5356 & 0.4470 & 0.2246 \\
eTour\_en & TF-IDF & 0.4102 & 0.6807 & 0.4753 & 0.3280 \\
eTour\_en & BM25 & 0.3074 & 0.5824 & 0.3704 & 0.2642 \\
RTMX-module-valid & TF-IDF & 0.7894 & 0.7894 & 0.8428 & 0.2500 \\
RTMX-module-valid & BM25 & 0.7606 & 0.7606 & 0.8194 & 0.2500 \\
SMOS & MiniLM & 0.2859 & 0.4568 & 0.2689 & 0.1984 \\
SMOS & Hybrid & 0.3026 & 0.5676 & 0.3092 & 0.2194 \\
iTrust & MiniLM & 0.2061 & 0.3394 & 0.2658 & 0.1422 \\
iTrust & Hybrid & 0.3068 & 0.4727 & 0.3723 & 0.1856 \\
eTour\_en & MiniLM & 0.4492 & 0.7113 & 0.5011 & 0.3508 \\
eTour\_en & Hybrid & 0.4406 & 0.7256 & 0.5280 & 0.3804 \\
RTMX-module-valid & MiniLM & 0.6795 & 0.6795 & 0.7592 & 0.2500 \\
RTMX-module-valid & Hybrid & 0.7515 & 0.7515 & 0.8138 & 0.2500 \\
SMOS & Hybrid + KG & 0.2965 & 0.5419 & 0.2953 & 0.2147 \\
iTrust & Hybrid + KG & 0.2958 & 0.4318 & 0.3677 & 0.1946 \\
eTour\_en & Hybrid + KG & 0.4744 & 0.7471 & 0.5415 & 0.3713 \\
SMOS & Hybrid + RAG & 0.3061 & 0.5875 & 0.3145 & 0.2194 \\
iTrust & Hybrid + RAG & 0.3043 & 0.4648 & 0.3691 & 0.1856 \\
eTour\_en & Hybrid + RAG & 0.4427 & 0.7232 & 0.5293 & 0.3804 \\
RTMX-module-valid & Hybrid + RAG & 0.6985 & 0.6985 & 0.7728 & 0.2500 \\
\end{longtable}

The eTour graph variant has MAP 0.4744 compared with hybrid 0.4406, but
its micro F1@10 is 0.3713 compared with hybrid 0.3804. This illustrates
why one favorable metric cannot summarize every retrieval objective. On
iTrust, BM25 has MAP 0.3787, higher than the reported semantic, hybrid,
graph and RAG variants. MiniLM alone leads hybrid on eTour MAP but not
its fixed-cutoff F1.

\begin{figure}
\centering
\includegraphics{../reports/results/traceability-map.png}
\caption{Measured MAP across the four datasets. RTMX has only seven
candidates; its scale must be interpreted separately.}
\end{figure}

\begin{figure}
\centering
\includegraphics{../reports/results/etour-precision-recall-f1.png}
\caption{eTour micro precision, recall and F1 at cutoff ten. These are
distinct from macro per-query recall.}
\end{figure}

\hypertarget{rag}{%
\section{RAG}\label{rag}}

\begin{longtable}[]{@{}
  >{\raggedright\arraybackslash}p{(\columnwidth - 8\tabcolsep) * \real{0.2000}}
  >{\raggedright\arraybackslash}p{(\columnwidth - 8\tabcolsep) * \real{0.2000}}
  >{\raggedright\arraybackslash}p{(\columnwidth - 8\tabcolsep) * \real{0.2000}}
  >{\raggedright\arraybackslash}p{(\columnwidth - 8\tabcolsep) * \real{0.2000}}
  >{\raggedright\arraybackslash}p{(\columnwidth - 8\tabcolsep) * \real{0.2000}}@{}}
\toprule\noalign{}
\begin{minipage}[b]{\linewidth}\raggedright
Dataset
\end{minipage} & \begin{minipage}[b]{\linewidth}\raggedright
Queries
\end{minipage} & \begin{minipage}[b]{\linewidth}\raggedright
Changed ranks
\end{minipage} & \begin{minipage}[b]{\linewidth}\raggedright
Uncertain judgments
\end{minipage} & \begin{minipage}[b]{\linewidth}\raggedright
Mean seconds
\end{minipage} \\
\midrule\noalign{}
\endhead
\bottomrule\noalign{}
\endlastfoot
SMOS & 67 & 17 & 217 & 70.82 \\
iTrust & 105 & 60 & 202 & 42.23 \\
eTour\_en & 57 & 33 & 121 & 37.93 \\
RTMX-module-valid & 22 & 19 & 37 & 67.28 \\
\end{longtable}

The RAG reranker slightly raises MAP on SMOS and eTour, lowers it on
iTrust, and lowers RTMX MAP from 0.7515 to 0.6985. The eTour result is
mixed: MAP and nDCG rise slightly while MRR and macro Recall@1 decline.
No statistical significance or general superiority follows. All 251
outputs pass the independent full-schema check, but that is not a
semantic grounding score.

\hypertarget{change-investigation}{%
\section{Change investigation}\label{change-investigation}}

All twelve authored scenarios contain their expected status somewhere in
the deterministic findings. The known affected sets contain 18 endpoints
in total, of which 14 are recovered and 4 missed. Predictions outside
these partial sets cannot be counted as verified false positives. These
counts are regression evidence under the authored contracts, not
independent impact precision/F1.

The recorded API-change demonstration changes the backend route prefix
from \texttt{/items} to \texttt{/catalog}. The generated frontend client
retains \texttt{/api/v1/items/\{id\}}. The executed workflow records
twelve tool stages, eleven changed artifact records, 33 impact
candidates and 24 verified excerpts. Its fifteen rule findings concern a
shared cause. The local model returns UNCERTAIN.

SAFA compares artifact-tree versions containing 240 versus 360 nodes and
234 versus 355 edges. The saved comparison identifies 205 added, 85
deleted and nine modified nodes, plus 207 added and 86 deleted edges.
This is temporal artifact analysis without independent consistency
labels.

\hypertarget{ablation-study}{%
\chapter{Ablation Study}\label{ablation-study}}

\begin{longtable}[]{@{}
  >{\raggedright\arraybackslash}p{(\columnwidth - 10\tabcolsep) * \real{0.1667}}
  >{\raggedright\arraybackslash}p{(\columnwidth - 10\tabcolsep) * \real{0.1667}}
  >{\raggedright\arraybackslash}p{(\columnwidth - 10\tabcolsep) * \real{0.1667}}
  >{\raggedright\arraybackslash}p{(\columnwidth - 10\tabcolsep) * \real{0.1667}}
  >{\raggedright\arraybackslash}p{(\columnwidth - 10\tabcolsep) * \real{0.1667}}
  >{\raggedright\arraybackslash}p{(\columnwidth - 10\tabcolsep) * \real{0.1667}}@{}}
\toprule\noalign{}
\begin{minipage}[b]{\linewidth}\raggedright
Variant
\end{minipage} & \begin{minipage}[b]{\linewidth}\raggedright
Calls
\end{minipage} & \begin{minipage}[b]{\linewidth}\raggedright
Agreement
\end{minipage} & \begin{minipage}[b]{\linewidth}\raggedright
Uncertainty
\end{minipage} & \begin{minipage}[b]{\linewidth}\raggedright
Citation validity
\end{minipage} & \begin{minipage}[b]{\linewidth}\raggedright
Known endpoint recall
\end{minipage} \\
\midrule\noalign{}
\endhead
\bottomrule\noalign{}
\endlastfoot
E & 12 & 0.250 & 1.000 & 1.000 & 0.000 \\
F & 12 & 0.250 & 1.000 & 0.917 & 0.000 \\
G & 12 & 0.250 & 1.000 & 0.917 & 0.200 \\
H & 12 & 0.250 & 1.000 & 0.917 & 0.320 \\
I & 12 & 0.250 & 1.000 & 0.833 & 0.720 \\
\end{longtable}

All five variants return UNCERTAIN for every case. Their 25\% agreement
with authored labels equals the always-UNCERTAIN comparator. Five
generated responses fail citation validation and become uncertain. More
elaborate evidence selection does not improve the observed final
consistency classification.

Known endpoint recall increases across some configurations, reaching
0.72 in I, but this measures coverage of the incomplete authored oracle.
It does not demonstrate correct reasoning over those endpoints. Exact
quote checking and endpoint coverage are useful process measurements
that must remain separate from semantic decision quality.

\begin{figure}
\centering
\includegraphics{../reports/ablation/comparison.png}
\caption{Recorded local-model ablation. All arms have the same final
status agreement and uncertainty rate.}
\end{figure}

Per-call times retain their original measurements, including reused
calls and unequal cold-load effects. Their mean values describe this
execution history; they are not a controlled speed comparison. The
project incurred zero model API charges for these local calls, but
compute and energy costs were not measured as zero.

\hypertarget{discussion}{%
\chapter{Discussion}\label{discussion}}

RQ1 has a dataset-dependent answer. No reported retrieval method
dominates every dataset and metric. Graph context can improve rank
quality on the eTour test split while reducing fixed-cutoff recall or
hurting other populations. RAG changes order without expanding the
retrieval set, which limits what its success can mean.

RQ2 is answered negatively for the tested local model and protocol.
Richer context and bounded investigation do not improve its consistency
agreement above always-UNCERTAIN. The experiment does not isolate
whether model capacity, excerpts, schema constraints, prompt
construction or label ambiguity causes the observed uncertainty.
Larger-model or alternative-context hypotheses require separate
experiments rather than retrospective assertions.

RQ3 is supported as an engineering execution claim. The pipeline
processes actual commits, persists typed relationships, retrieves source
excerpts and records tool outcomes. The API-route example provides
inspectable evidence of an obsolete client contract. The correctness and
usefulness of every affected candidate remain unproven.

RQ4 separates three validation layers: the source quote exists at the
recorded version; the output conforms to a schema and uses supplied
identifiers; and the conclusion is semantically justified. The first two
are checked. The third requires independent human annotation or another
validated assessment procedure. The full central claim of superior
evidence-grounded consistency decisions is therefore not established.

\hypertarget{limitations}{%
\chapter{Limitations}\label{limitations}}

The benchmark populations are small and heterogeneous. Unlinked
requirements are excluded. Public benchmark contamination in model
pretraining is unknown. The controlled changes are correlated, and their
labels and baseline trace links are authored. Natural software evolution
and independent adjudication are needed for broader generalization.

Static parsing cannot reliably resolve reflection, dynamic dispatch,
arbitrary route construction or every language version. The graph
includes heuristic relationships, whose confidence values are not
calibrated probabilities. Prefix excerpts and chunk limits can omit
relevant code. The top-five RAG policy cannot recover missing candidates
and has no recall gain at cutoffs five or ten by design.

The recorded runtime includes interruptions and cold-load effects. Saved
process RSS is a diagnostic snapshot rather than measured peak system
memory. Browser validation is unverified because the available browser
archive endpoint returned HTML instead of a usable archive. Windows
execution and public multiuser deployment have not been tested.

Reference checking covers selected primary metadata, abstracts and model
cards; broader full-text synthesis is still required for a comprehensive
survey. Institutional author details, final formatting, plagiarism
review, publication policy and venue indexing remain unresolved. The
draft should not be represented as accepted, submitted or independently
reviewed.

\hypertarget{conclusion}{%
\chapter{Conclusion}\label{conclusion}}

VIGILANT provides an executed and recoverable research implementation
connecting repository changes, canonical artifacts, structural graphs,
trace candidates and source-verifiable investigation records. The
completed evaluation contains 251 RAG generations, 60
consistency-ablation calls and twelve controlled change scenarios, with
raw evidence and independent arithmetic checks preserved.

The results support a narrower conclusion than the original aspiration.
Retrieval effects are mixed and the tested local model does not improve
consistency classification. These findings identify concrete limits of
the current configuration while preserving a platform for controlled
follow-up work. A successful execution trace and valid citations are
insufficient to establish semantic reliability.

\hypertarget{future-work}{%
\chapter{Future Work}\label{future-work}}

The highest-priority extension is independent annotation of naturally
occurring changes, including complete affected-artifact sets, trace
links and consistency labels. Reviewers should annotate evidence
sufficiency separately from final status, remain blind to model variant
where feasible, and resolve disagreements through a documented
procedure. Existing review packets are preparation materials, not
completed annotations.

Future model studies can vary chunk selection, candidate count, model
capacity and verification policy under a new frozen protocol. Candidate
expansion and order-only reranking should be evaluated separately.
Validation-based selection must precede final test scoring, and current
negative runs should remain preserved.

Additional engineering work includes dynamic relationship support,
broader parser coverage, browser/accessibility checks, user-machine
verification and appropriately secured deployment if ever required. The
thesis and paper drafts need author review, verified institutional
requirements and further literature coverage before final submission.

\hypertarget{references}{%
\chapter{References}\label{references}}

\begin{enumerate}
\def\labelenumi{\arabic{enumi}.}
\item
  S. Robertson and H. Zaragoza. \emph{The Probabilistic Relevance
  Framework: BM25 and Beyond}. 2009. DOI:
  \href{https://doi.org/10.1561/1500000019}{10.1561/1500000019}.
\item
  G. V. Cormack, C. L. A. Clarke and S. Buettcher. \emph{Reciprocal rank
  fusion outperforms condorcet and individual rank learning methods}.
  SIGIR 2009, pp.~758-759. DOI:
  \href{https://doi.org/10.1145/1571941.1572114}{10.1145/1571941.1572114}.
\item
  N. Reimers and I. Gurevych. \emph{Sentence-BERT: Sentence Embeddings
  using Siamese BERT-Networks}. EMNLP-IJCNLP 2019, pp.~3982-3992. DOI:
  \href{https://aclanthology.org/D19-1410/}{10.18653/v1/D19-1410}.
\item
  P. Lewis et al.~\emph{Retrieval-Augmented Generation for
  Knowledge-Intensive NLP Tasks}. arXiv:2005.11401, 2020.
  \href{https://arxiv.org/abs/2005.11401}{Primary record}.
\item
  D. Fuchss, T. Hey, J. Keim, H. Liu, N. Ewald, T. Thirolf and A.
  Koziolek. \emph{LiSSA: Toward Generic Traceability Link Recovery
  Through Retrieval-Augmented Generation}. ICSE 2025, pp.~1396-1408.
  DOI: \href{https://ardoco.de/c/icse25}{10.1109/ICSE55347.2025.00186}.
\item
  J. Keim, S. Corallo, D. Fuchss, T. Hey, T. Telge and A. Koziolek.
  \emph{Recovering Trace Links Between Software Documentation And Code}.
  ICSE 2024, article 215. DOI:
  \href{https://ardoco.de/c/icse24}{10.1145/3597503.3639130}.
\item
  J. Keim, S. Corallo, D. Fuchss and A. Koziolek. \emph{Detecting
  Inconsistencies in Software Architecture Documentation Using
  Traceability Link Recovery}. ICSA 2023, pp.~141-152. DOI:
  \href{https://ardoco.de/c/icsa23}{10.1109/ICSA56044.2023.00021}.
\item
  D. Edge et al.~\emph{From Local to Global: A Graph RAG Approach to
  Query-Focused Summarization}. arXiv:2404.16130, 2024; viewed revision
  v2, 2025. \href{https://arxiv.org/abs/2404.16130}{Primary record}.
\item
  J. Yang, C. E. Jimenez, A. Wettig, K. Lieret, S. Yao, K. Narasimhan
  and O. Press. \emph{SWE-agent: Agent-Computer Interfaces Enable
  Automated Software Engineering}. arXiv:2405.15793, 2024.
  \href{https://arxiv.org/abs/2405.15793}{Primary record}.
\item
  Sentence Transformers. \emph{all-MiniLM-L6-v2 model card}.
  \href{https://huggingface.co/sentence-transformers/all-MiniLM-L6-v2}{Model
  publisher}. Accessed 7 October 2026. Experimental revision recorded in
  \texttt{data/models.json}.
\item
  Qwen. \emph{Qwen2.5-1.5B-Instruct model card}.
  \href{https://huggingface.co/Qwen/Qwen2.5-1.5B-Instruct}{Model
  publisher}. Accessed 7 October 2026. Experimental revision recorded in
  \texttt{data/models.json}.
\end{enumerate}

\hypertarget{appendix-reproduction-and-evidence-index}{%
\chapter{Appendix: Reproduction and Evidence
Index}\label{appendix-reproduction-and-evidence-index}}

The unified project package contains the source, raw/prepared datasets,
Git histories, controlled commits, database snapshots, knowledge graph,
immutable experiments and derived outputs. Installed packages and model
weights are recreated from pinned manifests. Windows setup targets the
requested folder without asserting that this hosted session wrote there.

\begin{itemize}
\item
  lexical:
  \texttt{experiments/retrieval/20261006T020645-ee4c5e32/results.json}
\item
  semantic:
  \texttt{experiments/retrieval/20261006T021326-432d54b3/results.json}
\item
  graph:
  \texttt{experiments/retrieval/20261006T142659-bc349a3f/results.json}
\item
  ranker:
  \texttt{experiments/ranker/20261006T142709-dea10c63/results.json}
\item
  changebench:
  \texttt{experiments/changebench/20261006T143016-f81c6507/results.json}
\item
  ablation:
  \texttt{experiments/ablation/20261006T194639-30694b4d/results.json}
\item
  pipeline:
  \texttt{experiments/pipeline/20261006T195007-e56e6b8d/results.json}
\item
  safa: \texttt{experiments/safa/20261006T142403-c681a243/results.json}
\item
  rag\_SMOS:
  \texttt{experiments/retrieval/20261007T104905-4064b602/results.json}
\item
  rag\_iTrust:
  \texttt{experiments/retrieval/20261007T163255-0dfbb833/results.json}
\item
  rag\_eTour\_en:
  \texttt{experiments/retrieval/20261007T171402-64fa98d8/results.json}
\item
  rag\_RTMX-module-valid:
  \texttt{experiments/retrieval/20261007T210208-3f0803cc/results.json}
\end{itemize}

To inspect current progress, run
\texttt{python\ scripts/continue\_project.py\ -\/-status}. To reproduce
derived tables without model inference, run the two audit scripts,
\texttt{scripts/summarize\_rag.py} and \texttt{scripts/evaluate.py}. To
reopen the recorded demo, start the local API and inspect the report for
baseline \texttt{1762adac607a1b29cfc4da129557780beea71616} and changed
commit \texttt{892a3d5213e950be40467f4a0c8cfb33d12e8132}.

The automated continuation mechanism executes existing scripts and
preserves checkpoints. It does not conduct independent human review or
bypass account limits. Read the active lease, ledger and resume record
before any new run. Keep completed calls and experiments unchanged.

\end{document}
